\chapter{Évaluation et expérience reproductible}
\label{ch:evaluation}

\section{Matrice de confusion et mesures}

\begin{lecture}
L'exactitude compte la proportion totale de décisions correctes. La matrice de
confusion conserve davantage d'information : ses lignes indiquent les maisons
réelles et ses colonnes les maisons prédites. Elle montre ainsi quelles classes
sont bien reconnues et lesquelles sont confondues.
\end{lecture}

Pour un test de taille $m_{\mathrm{te}}$, l'exactitude est
\[
  \operatorname{accuracy}=\frac1{m_{\mathrm{te}}}
  \sum_{i\in\mathcal I_{\mathrm{te}}}\ind[\widehat c_i=c_i].
\]
Cette quantité est l'estimateur empirique de la probabilité de classification
correcte sur un échantillon test indépendant \cite[chap.~7]{hastie2009}.
En français statistique, \emph{exactitude} évite la confusion avec la
\emph{précision} au sens anglophone de $\mathrm{precision}=TP/(TP+FP)$.

L'évaluation reçoit deux vecteurs de même longueur : les classes réelles
$(c_i)$ et les classes prédites $(\widehat c_i)$ du test. Chaque paire
$(c_i,\widehat c_i)$ incrémente une case d'une matrice $K\times K$. L'exactitude
est calculée depuis la diagonale ; les mesures par classe utilisent ensuite les
sommes des lignes et des colonnes. Aucun réentraînement n'intervient dans ces
comptages.

La matrice de confusion $M$, définie par
$M_{ab}=\#\{i:c_i=a,\widehat c_i=b\}$, révèle quelles maisons sont
confondues. Pour chaque classe, on complète par rappel, précision et score $F_1$.
La moyenne macro donne le même poids à chaque maison ; la moyenne pondérée tient
compte de leurs effectifs. Une valeur unique ne remplace pas ces diagnostics.

Le score $F_1$ d'une classe est la moyenne harmonique de sa précision $P$ et de
son rappel $R$ :
\[
  F_1=\frac{2PR}{P+R}.
\]
Il n'est élevé que si les deux quantités le sont. Avec $P=0{,}80$ et $R=0{,}50$,
on obtient $F_1\approx0{,}615$, et non leur moyenne arithmétique $0{,}65$.

Supposons ensuite deux classes : la première contient $900$ observations et a un
$F_1$ de $0{,}95$ ; la seconde en contient $100$ et a un $F_1$ de $0{,}40$. La
moyenne macro vaut $(0{,}95+0{,}40)/2=0{,}675$. La moyenne pondérée vaut
$(900\times0{,}95+100\times0{,}40)/1000=0{,}895$. Le second nombre paraît bien
meilleur parce que la classe majoritaire domine le calcul. Rapporter les deux
indique si l'on privilégie la performance par individu ou l'équilibre entre
classes.

\begin{table}[htbp]
\centering
\setlength{\tabcolsep}{9pt}
\begin{tabular}{lrrrr|r}
\toprule
réelle $\backslash$ prédite & G & H & R & S & total \\
\midrule
G & \cellcolor{accent!35}$42$ & $3$ & $2$ & $3$ & $50$ \\
H & $4$ & \cellcolor{accent!35}$35$ & $6$ & $5$ & $50$ \\
R & $1$ & $4$ & \cellcolor{accent!35}$43$ & $2$ & $50$ \\
S & $5$ & $2$ & $3$ & \cellcolor{accent!35}$40$ & $50$ \\
\midrule
total prédit & $52$ & $44$ & $54$ & $50$ & $200$ \\
\bottomrule
\end{tabular}
\caption{Matrice de confusion fictive. Les bonnes décisions occupent la
diagonale. L'exactitude vaut $(42+35+43+40)/200=0{,}80$, mais la ligne Hufflepuff
révèle un rappel inférieur aux autres classes.}
\label{tab:confusion-example}
\end{table}

Dans cet exemple, le rappel de Hufflepuff vaut $35/50=0{,}70$, tandis que sa
précision vaut $35/44\approx0{,}80$. La première quantité répond à « quelle part
des vrais Hufflepuff est retrouvée ? » ; la seconde à « parmi les élèves prédits
Hufflepuff, quelle part l'est réellement ? ». Ces dénominateurs différents
expliquent pourquoi les deux mesures ne sont pas interchangeables.

\begin{figure}[htbp]
\centering
\begin{tikzpicture}[x=.12cm,y=1cm]
  \draw[-{Latex},thick] (0,0)--(105,0) node[right]{\%};
  \foreach \x in {0,20,40,60,80,100}{
    \draw (\x,.08)--(\x,-.08) node[below]{\x};
  }
  \foreach \yy/\name/\rec/\prec in
    {4/Gryffindor/84/80.8,3/Hufflepuff/70/79.5,2/Ravenclaw/86/79.6,1/Slytherin/80/80}{
    \node[anchor=east] at (-2,\yy){\name};
    \fill[accent!75] (0,\yy+.05) rectangle (\rec,\yy+.35);
    \fill[warning!70] (0,\yy-.35) rectangle (\prec,\yy-.05);
  }
  \node[accent,anchor=west] at (5,5){rappel};
  \node[warning,anchor=west] at (25,5){précision};
\end{tikzpicture}
\caption{Rappel et précision par classe calculés à partir de la
table~\ref{tab:confusion-example}. Une exactitude globale de $80\,\%$ masque
notamment le rappel plus faible de Hufflepuff.}
\label{fig:metrics-comparison}
\end{figure}

\begin{procedure}{Construire les mesures de classification}
\textbf{Entrées :} $m_{\mathrm{te}}$ classes réelles et autant de classes
prédites, codées dans le même ordre de $K$ libellés. \textbf{Sorties :} une
matrice de confusion $M$, l'exactitude et les mesures par classe.
\begin{enumerate}
  \item initialiser $M\in\mathbb{N}^{K\times K}$ à zéro ; les lignes représentent le
        réel et les colonnes la prédiction ;
  \item pour chaque paire $(c_i,\widehat c_i)$, incrémenter $M_{c_i,\widehat c_i}$ ;
  \item vérifier que la somme de toutes les cases vaut $m_{\mathrm{te}}$ ;
  \item calculer l'exactitude comme somme diagonale divisée par $m_{\mathrm{te}}$ ;
  \item pour la classe $k$, calculer rappel = diagonale/somme de la ligne et
        précision = diagonale/somme de la colonne, puis $F_1$ ;
  \item signaler tout dénominateur nul et rapporter les valeurs par classe.
\end{enumerate}
\end{procedure}

\section{Protocole minimal d'évaluation}

\begin{lecture}
Un score isolé est impossible à interpréter sans connaître la manière dont il a
été obtenu. Le protocole enregistre les données utilisées, leur préparation, les
paramètres d'entraînement et les mesures finales afin que la même expérience
puisse être reproduite et contrôlée.
\end{lecture}

Une expérience exploitable doit conserver :
\begin{enumerate}
  \item les variables retenues et la politique de valeurs manquantes ;
  \item la partition, sa stratification et sa graine ;
  \item les statistiques de préparation apprises sur le train ;
  \item $\alpha$, $\eps$, le nombre maximal d'itérations et, le cas échéant,
        $\lambda$ ;
  \item pour chaque classe, le coût initial, le coût final, le nombre
        d'itérations et la norme finale du gradient ;
  \item sur le test : matrice de confusion, exactitude et mesures par classe.
\end{enumerate}

Deux contrôles simples donnent du sens au résultat. Le modèle doit dépasser une
référence majoritaire qui prédit toujours la maison la plus fréquente. Puis une
permutation aléatoire des étiquettes doit ramener la performance vers le hasard ;
une performance anormalement élevée après permutation signale souvent une fuite
de données ou une erreur de protocole.

Exemple : si Gryffindor représente $35\,\%$ du test, prédire Gryffindor pour tout
le monde donne déjà $35\,\%$ d'exactitude. Un modèle obtenant $38\,\%$ ne doit
donc pas être comparé seulement au hasard uniforme de $25\,\%$, mais aussi à
cette référence majoritaire. Pour le contrôle par permutation, on mélange les
maisons entre les lignes avant l'entraînement tout en conservant les notes. Le
lien entre caractéristiques et réponse est alors détruit. Sur quatre classes
équilibrées, une exactitude proche de $25\,\%$ est attendue ; obtenir encore
$80\,\%$ indiquerait presque certainement que la réponse ou une information qui
la révèle s'est introduite dans les variables ou dans la préparation.

\begin{procedure}{Exécuter une évaluation contrôlée}
\textbf{Entrées :} modèle figé, statistiques de préparation, indices de test et
classes réelles. \textbf{Sortie :} un compte rendu reproductible, sans
modification du modèle.
\begin{enumerate}
  \item sélectionner les lignes désignées par $\mathcal I_{\mathrm{te}}$ et leur
        appliquer les statistiques d'apprentissage déjà conservées ;
  \item produire exactement une prédiction par ligne, sans mise à jour des paramètres ;
  \item vérifier l'égalité des longueurs entre prédictions et classes réelles ;
  \item construire matrice de confusion, exactitude et mesures par classe ;
  \item comparer l'exactitude à celle de la classe majoritaire et exécuter
        séparément le contrôle par permutation sur une expérience d'apprentissage ;
  \item enregistrer résultats, paramètres, hyperparamètres, graine et indices.
\end{enumerate}
\end{procedure}

\section{Validation croisée : plis et répétitions}

\begin{lecture}
Une seule partition peut avantager ou désavantager un modèle par hasard. La
validation croisée fait tourner le rôle de validation entre plusieurs blocs :
chaque observation est évaluée une fois tout en restant hors de l'apprentissage
correspondant. Répéter l'ensemble du découpage mesure la sensibilité à cette
répartition aléatoire.
\end{lecture}

Un \emph{pli} désigne un bloc d'observations. Dans une validation croisée
à cinq plis, on découpe les données disponibles pour la sélection du modèle en
cinq blocs disjoints $B_1,\ldots,B_5$. On entraîne cinq fois le modèle : à chaque
fois, quatre blocs servent à l'ajustement et le bloc restant sert à la validation.

\begin{figure}[htbp]
\centering
\begingroup
\setlength{\tabcolsep}{4pt}
\renewcommand{\arraystretch}{1.35}
\begin{tabular}{r*{5}{>{\centering\arraybackslash}p{13mm}}}
& $B_1$ & $B_2$ & $B_3$ & $B_4$ & $B_5$ \\
validation sur $B_1$ & \cellcolor{warning!70}V & \cellcolor{accent!55}A & \cellcolor{accent!55}A & \cellcolor{accent!55}A & \cellcolor{accent!55}A \\
validation sur $B_2$ & \cellcolor{accent!55}A & \cellcolor{warning!70}V & \cellcolor{accent!55}A & \cellcolor{accent!55}A & \cellcolor{accent!55}A \\
validation sur $B_3$ & \cellcolor{accent!55}A & \cellcolor{accent!55}A & \cellcolor{warning!70}V & \cellcolor{accent!55}A & \cellcolor{accent!55}A \\
validation sur $B_4$ & \cellcolor{accent!55}A & \cellcolor{accent!55}A & \cellcolor{accent!55}A & \cellcolor{warning!70}V & \cellcolor{accent!55}A \\
validation sur $B_5$ & \cellcolor{accent!55}A & \cellcolor{accent!55}A & \cellcolor{accent!55}A & \cellcolor{accent!55}A & \cellcolor{warning!70}V \\
\end{tabular}

\vspace{.6em}
\textcolor{accent}{A : apprentissage}\qquad
\textcolor{warning}{V : validation}
\endgroup
\caption{Validation croisée à cinq plis. Chaque observation sert quatre fois à
l'apprentissage et exactement une fois à la validation. Le test final reste à
l'écart de l'ensemble de ce mécanisme.}
\label{fig:five-fold}
\end{figure}

Supposons que les exactitudes de validation des cinq tours soient
$0{,}79$, $0{,}82$, $0{,}77$, $0{,}81$ et $0{,}80$. Leur moyenne vaut
$0{,}798$. Leur dispersion montre aussi que le résultat dépend quelque peu des
observations placées dans le bloc de validation ; annoncer seulement « environ
$0{,}80$ » masque cette information.

Une \emph{répétition} consiste à refaire tout le découpage en cinq blocs avec une
autre permutation aléatoire, puis à recommencer les cinq entraînements. Les
observations ne changent pas ; seule leur affectation aux blocs change. Deux
répétitions produisent ici dix mesures :

\begin{table}[htbp]
\centering
\begin{tabular}{lccccc|c}
\toprule
& pli 1 & pli 2 & pli 3 & pli 4 & pli 5 & moyenne \\
\midrule
répétition 1 & $0{,}79$ & $0{,}82$ & $0{,}77$ & $0{,}81$ & $0{,}80$ & $0{,}798$ \\
répétition 2 & $0{,}81$ & $0{,}78$ & $0{,}80$ & $0{,}76$ & $0{,}83$ & $0{,}796$ \\
\midrule
\multicolumn{6}{r}{moyenne des dix validations} & $0{,}797$ \\
\bottomrule
\end{tabular}
\caption{Exemple fictif de validation croisée répétée. Les dix nombres ne sont
pas dix évaluations finales indépendantes : les ensembles d'apprentissage se
recouvrent fortement. Ils décrivent néanmoins la sensibilité au découpage.}
\label{tab:repeated-cv}
\end{table}

\begin{procedure}{Comparer des configurations par validation croisée}
\textbf{Entrées :} les données d'apprentissage, une liste de configurations et
un nombre de plis $q$. Un pli est une liste d'indices utilisée une fois pour la
validation. \textbf{Sortie :} une configuration choisie sans consulter le test.
\begin{enumerate}
  \item répartir les indices en $q$ blocs disjoints et enregistrer ces blocs ;
  \item pour une configuration donnée et un bloc $B_r$, ajuster sur tous les
        indices hors de $B_r$, puis mesurer sur $B_r$ ;
  \item répéter pour $r=1,\ldots,q$ : chaque ligne sert ainsi une fois à la validation ;
  \item calculer la moyenne et la dispersion des $q$ mesures ;
  \item recommencer avec exactement les mêmes blocs pour chaque configuration ;
  \item choisir selon une règle fixée, puis seulement exécuter le protocole final et le test.
\end{enumerate}
\end{procedure}

Pour comparer deux modèles A et B, on utilise exactement les mêmes blocs pour les
deux. Si, sur un pli donné, A obtient $0{,}81$ et B $0{,}79$, la différence est
$+0{,}02$. Répéter ce calcul pli par pli sépare mieux l'effet du modèle de la
difficulté particulière d'un bloc. Employer des découpages différents ajouterait
inutilement la variation des données à la comparaison.

\begin{vigilance}[Incertitude de l'évaluation]
Une exactitude test est elle-même une estimation soumise à la variabilité
d'échantillonnage. Sur $200$ observations, une différence de deux bonnes
prédictions ne représente qu'un point de pourcentage et ne suffit généralement
pas à conclure qu'un modèle est réellement supérieur. La comparaison de modèles
doit conserver les mêmes partitions. Les plis et répétitions de la section
précédente mesurent la sensibilité pendant la sélection ; ils ne transforment pas
le test final en certitude.
\end{vigilance}

Par exemple, A classe correctement $162$ élèves sur $200$ et B en classe $160$.
L'écart d'exactitude vaut $0{,}01$, mais il faut regarder les décisions appariées.
Si A corrige deux erreurs de B sans introduire aucune nouvelle erreur, l'écart est
cohérent sur ces deux cas. Si A et B se trompent chacun sur vingt élèves presque
tous différents, le même total cache des comportements très différents. La
matrice de confusion et la liste des désaccords sont alors plus informatives que
le seul écart d'un point.

\begin{resultat}[Résultats produits par l'évaluation]
L'évaluation finale produit le vecteur des classes prédites, la matrice de
confusion, l'exactitude globale et les mesures par classe. Les historiques de
coût et de norme du gradient décrivent la convergence de l'entraînement ; ils ne
remplacent pas les mesures calculées sur le test.
\end{resultat}
